\documentclass[a4paper,12pt]{article}

\usepackage[margin=0.5in]{geometry}
\usepackage{graphicx}
\usepackage{fontspec}
\usepackage{enumitem}
\usepackage{hyperref}
\usepackage{setspace}
\usepackage{xcolor}
\usepackage[most]{tcolorbox}
\usepackage{longtable}
\usepackage{booktabs}
\usepackage{minted}

\onehalfspacing

% For the specific font family
\newfontfamily\majormono{Major Mono Display}[
    Path = ./fonts/,
    Extension = .ttf
]

\setmainfont[
    Path = ./fonts/,
    Extension = .ttf,
    UprightFont = *-Regular,
    BoldFont = *-Bold,
]{Outfit}

\setlength{\parindent}{0pt}
\setlength{\parskip}{0.5em}

\hypersetup{
    colorlinks=true,
    linkcolor=blue,
    urlcolor=blue
}

\newtcolorbox{ghnote}{
    colback=gray!5,
    colframe=gray!40,
    boxrule=0.4pt,
    arc=4pt,
    left=6pt, right=6pt, top=6pt, bottom=6pt
}

\newtcbox{\keys}{
    on line,
    boxsep=2pt,
    left=2pt,
    right=2pt,
    top=1pt,
    bottom=1pt,
    colback=gray!10,
    colframe=gray!40,
    boxrule=0.4pt,
    arc=3pt,
    fontupper=\ttfamily\small
}

\begin{document}

\hspace{-0.25em}{\majormono\fontsize{30pt}{36pt}\selectfont lektra}

High-performance Document and Image viewer that prioritizes screen space and control.

\vspace{1em}
\hrule
\vspace{1em}

\section*{Important Links}

\begin{description}[labelwidth=2.2cm, labelsep=0.8em, align=right, leftmargin=0pt]
    \item[\textbf{Homepage}]       \url{https://dheerajshenoy.github.io/lektra}
    \item[\textbf{Configuration}]  \url{https://dheerajshenoy.github.io/lektra/configuration.html}
    \item[\textbf{Commands}]       \url{https://dheerajshenoy.github.io/lektra/commands.html}
    \item[\textbf{Code}]           \url{https://codeberg.org/lektra/lektra}
\end{description}

\begin{ghnote}
    \textbf{Note:} This guide assumes you have installed \textbf{LEKTRA} and have not changed the default keybindings.
\end{ghnote}

\tableofcontents
\newpage

% ---------------------------------------------------------------------------
\section{Getting Started}
% ---------------------------------------------------------------------------

\subsection{Opening a file}
You can open a document in multiple ways:

\begin{itemize}[leftmargin=*]
    \item \textbf{Command line:} \keys{lektra <path-to-file>}
    \item \textbf{File explorer:} Right-click a file and choose \textit{Open with lektra}
    \item \textbf{From within lektra:} Press \keys{o} to open the file picker
\end{itemize}

Supported formats include PDF, EPUB, XPS, CBZ, and common image formats. DjVu support is detected at runtime if \texttt{djvulibre} is installed.

\subsection{Tabs and Splits}

LEKTRA supports tabs and splits for viewing multiple documents simultaneously.

\begin{itemize}[leftmargin=*]
    \item \keys{Ctrl+t} opens a new tab
    \item \keys{Ctrl+w} closes the current tab
    \item \keys{gt} / \keys{gT} cycles to the next / previous tab
    \item \keys{Ctrl+\textbackslash} splits the view vertically
    \item \keys{Ctrl+-} splits the view horizontally
\end{itemize}

\subsection{Keyboard Navigation and Keybindings}

The core philosophy of \textbf{LEKTRA} is keyboard-driven navigation (mouse support is available too). Default keybindings follow VIM conventions.

\subsubsection{Scrolling and Movement}

\begin{itemize}[leftmargin=*]
    \item \textbf{Vertical:} \keys{j} scrolls down,\; \keys{k} scrolls up
    \item \textbf{Horizontal:} \keys{h} scrolls left,\; \keys{l} scrolls right
\end{itemize}

\subsubsection{Page Navigation}

\begin{itemize}[leftmargin=*]
    \item \keys{Shift+j} moves to the next page
    \item \keys{Shift+k} moves to the previous page
    \item \keys{g,g} jumps to the first page
    \item \keys{Shift+g} jumps to the last page
    \item \keys{Ctrl+g} go to a specific page (prompts for page number)
\end{itemize}

\subsubsection{Search}
\begin{itemize}[leftmargin=*]
    \item \keys{/} opens or focuses the search bar
    \item \keys{n} jumps to the next search result
    \item \keys{Shift+n} jumps to the previous search result
\end{itemize}

\subsubsection{Zoom and View Control}
\begin{itemize}[leftmargin=*]
    \item \keys{=} zooms in, \quad \keys{-} zooms out, \quad \keys{0} resets zoom
\end{itemize}

Fit modes:
\begin{itemize}[leftmargin=*]
    \item \keys{Ctrl+Shift+W} fit to width
    \item \keys{Ctrl+Shift+H} fit to height
    \item \keys{Ctrl+Shift+=} fit to window
    \item \keys{Ctrl+Shift+R} toggle auto-resize
\end{itemize}

\subsubsection{Outline and History}
\begin{itemize}[leftmargin=*]
    \item \keys{t} opens the document outline (table of contents) overlay
    \item \keys{Alt+Shift+H} opens the highlight annotation search overlay
    \item \keys{Ctrl+o} goes back to the previous location in history
\end{itemize}

\subsubsection{Annotations and Interaction Modes}
\begin{ghnote}
    \textbf{Note:} These keys toggle modes. Some modes require clicking or dragging on the page to take effect.
\end{ghnote}

\begin{itemize}[leftmargin=*]
    \item \keys{1} toggle text selection mode
    \item \keys{2} toggle text highlight mode
    \item \keys{3} toggle rectangle annotation mode
    \item \keys{4} toggle region selection mode
    \item \keys{5} toggle annotation popups
\end{itemize}

\subsubsection{Links and Actions}
\begin{itemize}[leftmargin=*]
    \item \keys{f} link-hint visit (keyboard link navigation)
    \item \keys{o} open file
    \item \keys{Ctrl+s} save current document
\end{itemize}

\subsubsection{Editing and UI}
\begin{itemize}[leftmargin=*]
    \item \keys{u} undo, \quad \keys{Ctrl+r} redo
    \item \keys{i} invert colors
    \item \keys{Ctrl+Shift+m} toggle menubar
    \item \keys{:} open command palette
\end{itemize}

\subsubsection{Rotation}
\begin{itemize}[leftmargin=*]
    \item \keys{<} rotate counter-clockwise, \quad \keys{>} rotate clockwise
\end{itemize}

% ---------------------------------------------------------------------------
\section{Configuration}
% ---------------------------------------------------------------------------

LEKTRA can be configured using \textbf{TOML} or \textbf{Lua}. Let us look at how to configure using TOML first. The configuration file should be at \keys{\$HOME/.config/lektra/config.toml}.

\subsection{Key sections}

\begin{description}[leftmargin=0pt]
    \item[\texttt{[window]}] Window size, fullscreen, menubar visibility, title format, accent color.
    \item[\texttt{[page]}] Page foreground and background colors.
    \item[\texttt{[statusbar]}] Visibility and which fields to display (page number, progress, file info, etc.).
    \item[\texttt{[split]}] Split behavior — focus follows mouse, dim inactive panes, etc.
    \item[\texttt{[portal]}] Picture-in-picture portal overlay settings.
    \item[\texttt{[annotations.*]}] Highlight, rectangle, and popup annotation appearance.
    \item[\texttt{[thumbnail\_panel]}] Thumbnail strip width, page numbers, and scrollbar visibility.
    \item[\texttt{[keymap]}] Override any keybinding by specifying \texttt{action = "Key"}.
    \item[\texttt{[mousemap]}] Mouse button bindings, similar to keymap.
\end{description}

\subsection{Example snippet}

\begin{tcolorbox}[colback=gray!5, colframe=gray!40, boxrule=0.4pt, arc=4pt,
    fontupper=\ttfamily\small]
\begin{minted}{toml}
[window]
menubar = false
accent = "#3daee9FF"

[statusbar]
visible = true
show_page_number = true
show_progress = true

[page]
bg = "#1e1e2e"
fg = "#cdd6f4"
\end{minted}
\end{tcolorbox}

For the full list of options see the \href{https://dheerajshenoy.github.io/lektra/configuration.html}{configuration reference} on the homepage.

% ---------------------------------------------------------------------------
\section{Commands}
% ---------------------------------------------------------------------------

Press \keys{:} to open the command palette and run any command by name. Commands cover opening files, navigation, zoom, annotations, sessions, and more.

The full command reference is on the \href{https://dheerajshenoy.github.io/lektra/commands.html}{commands page} of the homepage.

% ---------------------------------------------------------------------------
\section{Lua API}
% ---------------------------------------------------------------------------

LEKTRA embeds a Lua 5.x interpreter (build with \texttt{-DWITH\_LUA=ON}). All APIs live under the global \texttt{lektra} table.

\subsection{Configuration file}

Place your Lua init file at \keys{\$HOME/.config/lektra/init.lua}. It is sourced at startup after the TOML config is applied.

\subsection{Namespaces}

\begin{description}[leftmargin=0pt]
    \item[\texttt{lektra.view}] Document view helpers — navigation, zoom, rotation, layout, search, text extraction, outline, and per-view event listeners.
    \item[\texttt{lektra.opt}]  Read/write runtime options (fit mode, layout mode, DPR, statusbar fields, …).
    \item[\texttt{lektra.cmd}]  Execute any lektra command by name from Lua.
    \item[\texttt{lektra.keymap}] Bind keys to Lua functions at runtime.
    \item[\texttt{lektra.mousemap}] Bind mouse buttons to Lua functions.
    \item[\texttt{lektra.event}] Register global event listeners (document open/close, page change, …).
    \item[\texttt{lektra.ui}]   Show message bars, input prompts, and pickers from Lua.
    \item[\texttt{lektra.bookmark}] Programmatic bookmark management.
    \item[\texttt{lektra.tabs}] Tab management — open, close, focus, list tabs.
    \item[\texttt{lektra.utils}] Utility helpers (shell execute, clipboard, path helpers).
\end{description}

\subsection{Quick examples}

\textbf{Bind a key to a Lua function:}

\begin{tcolorbox}[colback=gray!5, colframe=gray!40, boxrule=0.4pt, arc=4pt,
    fontupper=\ttfamily\small]
\begin{minted}{lua}
lektra.keymap.map("n", "<C-e>", function()
local v = lektra.view.current()
lektra.ui.message("Page: " .. v:pageno())
end)
\end{minted}
\end{tcolorbox}

\textbf{React to a page change event:}

\begin{tcolorbox}[colback=gray!5, colframe=gray!40, boxrule=0.4pt, arc=4pt,
    fontupper=\ttfamily\small]
\begin{minted}{lua}
local v = lektra.view.current()
v:register("page_changed", function()
print("Now on page", v:pageno())
end)
\end{minted}
\end{tcolorbox}

\textbf{Open a file in a new tab:}

\begin{tcolorbox}[colback=gray!5, colframe=gray!40, boxrule=0.4pt, arc=4pt,
    fontupper=\ttfamily\small]
\begin{minted}{lua}
lektra.tabs.open("/path/to/document.pdf")
\end{minted}
\end{tcolorbox}

\subsection{Type stubs}

LSP-friendly type stubs for the Lua API are installed at
\keys{\$prefix/share/lektra/lua/} (Linux/macOS). Point your Lua language server at this directory for autocompletion.

\begin{tcolorbox}[colback=gray!5, colframe=gray!40, boxrule=0.4pt, arc=4pt,
    fontupper=\ttfamily\small]
\begin{minted}{json}
// .luarc.json example
{
    "workspace.library": ["/usr/share/lektra/lua"]
}
\end{minted}
\end{tcolorbox}

% For the complete API reference with all methods, types, and event names, see the \href{https://dheerajshenoy.github.io/lektra/lua-api.html}{Lua API} page on the homepage.

\end{document}

